\documentclass[11pt]{article}
\usepackage[margin=1in]{geometry}
\usepackage{microtype}
\usepackage{booktabs}
\usepackage{enumitem}
\usepackage[hidelinks]{hyperref}
\usepackage{amsmath}
\usepackage{seqsplit}
\usepackage{graphicx}
\emergencystretch=2em
\title{Observer-Scoped Evidence for Time-Sensitive Claims\\\large Part I: A Self-Audit of Gate Vacuity and a Model-Free Contract Feasibility Check}
\author{Author list and affiliations to be completed after external review}
\date{October 5, 2026}

\begin{document}
\maketitle

\begin{abstract}
An agent answering time-sensitive questions should distinguish observer-accessible evidence from hidden world state. Valid time, transaction time and source provenance are established concepts; this draft studies their executable representation and the instruments meant to verify it. A model-free self-audit of one system's evaluation gates, spanning 288 deterministic projection conditions and 216 reference-sensitive nested arms, found that a hidden-world perturbation was discarded before the compared artifacts were built, so the equality it fed could not fail; the noninterference claim resting on it is withdrawn. Generalising that failure produced a register of nine structurally distinct mechanisms by which a deterministic, model-free assertion can be unfalsifiable, one entry recursive: the report that no vacuous conjunct remained was itself an assertion no perturbation reached. Each is stated system-independently with a detection question a reader can run on their own gate. Against an external benchmark at a pinned commit, a paired metamorphic probe withdrew a second inference of ours: representation equivalence is not verdict-preserving for that scorer, whose word-boundary guard applies only inside an ASCII character class, so three of six candidate equivalence operators move the verdict in both directions; a benign-transformation class must be certified per operator. On the benchmark's published corpus, 45 identifier-style targets (44 prefixes and 1 suffix) are verbatim present yet unmatched, but recomputing per checkpoint the number of context-side verdicts that would change is zero, backed by a self-test fixture that must flip. The confirmatory-query gate reports FAIL with all 84 queries still placeholders. No semantic rating, external baseline, model-quality or general-method claim is made.
\end{abstract}

\paragraph{Keywords} observer-scoped evidence; time-scoped evidence validity; refresh sampling; fault-targeted differential audit; gate vacuity; reproducible evaluation

\section{Introduction}
Agents operating in dynamic environments must distinguish what an observer has sensed, received, or done from global simulator state, stale support, plans, and historical records. Temporal-memory benchmarks such as LongMemEval and point-in-time character tests such as TimeChara address related reasoning problems, product-grade memory layers already store bi-temporal facts, and recent living-world benchmarks study long-horizon interaction. Our narrower question concerns the evidence acquisition boundary: which time-stamped facts are admissible to one observer, when were they sampled, and can an inaccessible world change affect the resulting context?

This manuscript reports protocol feasibility and a protocol self-audit, not answer-model performance. Its artifacts support reproducible failure analysis and candidate adequacy requirements, without establishing a novel gate method. Timeline families share a frozen authoring procedure; they are not independent human-authored instances. We ask whether observer evidence, refresh sampling, and current/history projection can be represented and replayed while keeping evaluator-only labels outside the observer-visible projection.

\section{Related Work and Positioning}
\paragraph{Temporal validity is not our contribution.} The split between when a fact holds and when a system came to know it is the standard bi-temporal model of temporal databases, formalised in the taxonomy of Snodgrass and Ahn and in the TSQL2 language specification, extended to dependency theory for temporal databases by Jensen, Snodgrass and Soo (1996), and the same split appears in the W3C sensor ontology as \texttt{phenomenonTime} versus \texttt{resultTime}. SQL dialects already expose half-open as-of semantics of the shape our projector uses. Indeterminate validity, the ancestor of the \texttt{unknown} column we still do not emit, was named valid-time indeterminacy in 1993.

\paragraph{Agent memory already implements time-scoped facts.} Zep and its Graphiti engine maintain a bi-temporal model with four timestamps per edge and invalidate, rather than delete, superseded facts; the expiry it records is inferred by the system when it detects a temporal contradiction, not declared by the source at record time. A-TMA separates bank, retrieval and answer-time failures for ``ghost memory'' and exposes current, historical and transition labels to the question answerer. WorldLines evaluates embodied assistants with an observer-grounded, visibility-aware memory. A 2026 line of benchmarks measures stale-memory effects directly: STALE on implicitly conflicting observations, memory-pollution revocation, the flattening of the temporal shape of user facts during memory consolidation across three installed pipelines, and outdated-document interference in retrieval. TRACE governs memory validity in evolving multi-agent systems through authenticated returns and reports valid-information availability separately from invalid-information rejection, so validity governance is an active design space rather than an open problem (Xiong et al., 2026).

\paragraph{Observer scoping and as-of gold are not our contribution either.} Three 2026 works occupy the ground this manuscript's framing would otherwise appear to claim. An audience-bound persistent memory records the audience present when an item is written, widens it only by explicit object-specific grant, admits an item into a model attempt only when every current viewer belongs to one of its authorized audiences with unresolved viewers failing closed to public-only, and reports soundness and policy-completeness under explicit identity, provenance and complete-mediation assumptions, confirmed prospectively over ten thousand multi-party histories in which unscoped retrieval exposed a forbidden item in 82 percent of contexts (Liu, 2026). That is an observer-scoped, fail-closed admission contract with provenance and a frozen confirmation design, so nothing about scoping evidence to a viewer is contributed here. A separate study of cross-family interference states the same division of labour as its conclusion: certification determines whether an edit is supported, while retrieval scope determines where that evidence authorizes its use, and it measures the difference with an execution-grounded gate over twenty-seven paired randomized-order streams (Cheng et al., 2026). A benchmark for authorization-limited evidence in agentic systems is the closest existing suite to the evidence-scoping half of this contract; it evaluates agents on evidence they are not permitted to see rather than auditing the gate that decides permission (Tallam, 2026). On the evaluation side, a drift-aware benchmark already writes every change to an append-only ground-truth ledger, computes gold answers from that ledger rather than from the live stores, resolves ground truth at both timestamps, and labels an answer that was correct when retrieved but wrong when evaluated as a freshness error distinct from a reasoning error (Singh et al., 2026). As-of gold resolution and the freshness-versus-reasoning split are therefore published engineering, not a contribution of this paper; what remains local is that the validity interval is declared by the source at record time rather than inferred by the system afterwards. A companion safety line already uses the phrase temporal authorization for trajectory-level gating of deterministic pre-execution checks (Chitan, 2026), so this manuscript deliberately avoids that term and speaks of time-scoped evidence validity instead. Finally, misattribution of a memory-layer fault to the model has been named and measured directly, with a counterfactual composition test that identifies the causal entry where forensic baselines fail (Ahad et al., 2026); the target-alignment control in Section 4 is an instance of that concern specialised to evaluator-only labels, not a discovery of it.

\paragraph{Abstention and the third label.} Supporting / refuting / not-enough-information labelling is standard in fact verification, unanswerability evaluation for retrieval exists with named metrics, temporal-question-answering abstention has been studied explicitly, and abstention is measured as selective risk against coverage in the selective-prediction literature. Any ``say I do not know'' affordance in this protocol is an instance of that work, not a new mechanism.

\paragraph{Noninterference and gate adequacy are also not new.} Noninterference originates in the security-policy literature of the early 1980s; comparing two executions to reduce it to a safety problem is self-composition, and testing noninterference against a reference implementation is published work. An assertion that the target fault cannot refute, because of how the test itself is constructed, is vacuous; measuring whether a test suite can detect a fault is mutation analysis. In natural-language evaluation, the corresponding hygiene is a trivial-baseline check. Mutation testing has since been turned on the evaluators themselves: rubric generation guided by response defects (Yang et al., 2026), semantic judges probed with eleven mutation operators (Delcourt et al., 2026), benchmark rigor audited with seeded mutants and explicit stillborn, equivalent and unsupported verdict classes (Bhadra, 2026), and reproducibility safeguards measured outcome-blind with equivalent mutants excluded from the denominator (Shulepov, 2026). Those works already supply operator registries, fault-activation requirements and layered denominators, so none of that apparatus is contributed here. A parallel engineering line separates comparison admission from single-arm completion and rejects invalid comparisons by replay and mutation before any result is reported (Fu et al., 2026). Most directly, a deterministic mutation-testing protocol for task-scoped state oracles in software-agent benchmarks already injects six classes of harmful state faults and three schema-certified benign transformations, compares official scores against controlled reference rules, reports an applicable-mutation denominator with Wilson intervals, and explicitly declines a benchmark ranking when one suite contributes too few applicable mutations (Sun, 2026). That is the protocol this manuscript had planned; it is prior work and is not claimed here. In the agent setting it has likewise been shown that removing the anchor guards from a self-evolving metric collapses it into a vacuous always-pass detector, and that a task score cannot validate the evaluator that produced it (Zhang et al., 2026). A systematic review of eighty-three studies sorts LLM test oracles by where their authority comes from rather than by their form, and reports that oracle quality is most often judged by resemblance to a known oracle rather than by whether injected faults are caught (Mughal and Bilal, 2026); that sentence is the gap this manuscript works inside, and it is not our finding. Two further separations are already established and are adopted here rather than proposed: fault triggering is not fault detection, and adequacy criteria that reach a faulty path still leave detection rates near one to two percent because the generated oracle is wrong (Hamidi et al., 2026); and a nested sample is not an independent unit, so a rule-level battery reports its comparison unit as the rule or credential type rather than the number of generated strings (Mishra, 2026). Metamorphic testing is the standard name for asserting a relation between the outputs of related inputs without an explicit oracle, and it has been surveyed and measured at scale for language models (Cho et al., 2025). Deterministic source-grounded admission gates that keep unverifiable candidates in a review representation rather than silently dropping them, and that refuse to treat twenty-two passing conformance tests as evidence beyond the exercised path, are likewise published engineering (Girda and Groza, 2026). The same hazard has been reported for a different instrument: an LLM-based security-log evaluator whose parsing silently suppresses entries it was built to score, so that the ruler rather than the measured system is where the error lives (Garware and Zisad, 2026). The matcher guard in Section 4 is a distinct mechanism in a distinct instrument, and that work is cited as evidence that instrument-side suppression is a recognised failure class, not as a precedent for this paper's specific finding. That evaluation reporting must carry its methodology and not only its scores is the argument of the holistic agent-evaluation infrastructure line (Kapoor et al., 2025); it is the reason this manuscript publishes a denominator, a detection-power statement and a vacuity disclosure beside every count it reports. One further item belongs here rather than in the contribution paragraph, because it is prior work and not a gap. An audit of an evidence-bundle verifier publishes a named definition: ``A verifier exhibits a vacuous pass when it returns acceptance while an intended binding, comparison, or coverage obligation was absent, bypassed, or semantically unexamined'' (Hill, 2026, Definition 1), together with the common ancestor it identifies for the class---an absence-assertion carrying no liveness proof---and a recorded instance of the recursive form this manuscript's register also carries. The vocabulary and the class are therefore published. What is contributed here is narrower: nine structurally distinct mechanisms enumerated inside one deterministic, model-free evaluation-gate suite, each bound to a retained artifact and a script digest, and the per-mechanism treatment in the Anatomy section. A prior round of review reported that the recursive entry had no occupier; that report searched with this manuscript's own vocabulary, and one query in the neighbouring vocabulary found this paper, so the retrieval procedure and not the literature was where the gap was.

\paragraph{What remains as a bounded contribution.} The artifact separates observer evidence, source-attested validity, evaluator-only labels, and gate failure records; that separation is bookkeeping and is not offered as a claim. Three residual points survive the preceding paragraphs, and each is narrow. First, and strongest, the gate-status table classifies nine structurally distinct reasons why a deterministic, model-free assertion could not fail. Published mutation audits register verdict classes for mutants---killed, survived, stillborn, equivalent, unsupported---which is an axis about the mutant; a reviewed published oracle-side property names and measures a single mechanism of the same family as M5 rather than enumerating these nine; Hill's verifier-side audit separately enumerates five fail-open and vacuous-pass classes but does not supply this nine-entry, per-mechanism treatment (Canedo, 2026; Hill, 2026). This register is an axis about the assertion, and those axes do not coincide. The English version of this paragraph carried the unnarrowed universal for one revision longer than its Chinese counterpart did, which is the worse place for it to have survived. One entry is recursive: the claim that no vacuous conjunct remained was itself an assertion that no input could refute, and it survived one re-issue before a second found it. Second, a gate whose perturbation set yields no sensitive arm reports \texttt{INCONCLUSIVE} rather than \texttt{PASS}, which is a verdict-semantics choice and not a detection technique; it is currently implemented in only one of the two gates that need it. Third, a benign-transformation class must be certified per operator against the scorer that consumes it: Section 4 shows three of six candidate representation-equivalence operators changing the verdict of an external matcher, in both directions. The hazard behind the GateMem fixture---a mutation of a secret must co-transform every evaluator-only target that names it, and a fixture whose labels no longer match its data is not a legal mutant---is handled constructively elsewhere, where gold labels are derived from the mutation itself so that co-transformation cannot be forgotten (de Oliveira et al., 2026), and as equivalent-mutant exclusion; what we found no published counterpart for is stating it as an explicit mutation-operator contract together with a demonstrated counterexample, and we claim that conjunction only. We therefore treat the current paper as a model-free self-audit of one system's own gates together with a contract feasibility check, not as a new noninterference method, not as a demonstrated advantage of source-attested validity, and not as a cross-system measurement. A broader measurement claim requires data-plane mutations, a falsification baseline, and multiple external harnesses.

\section{Methods}
\subsection{Artifact planes}
Each timeline family contains a verifier-only source registry, an observer evidence stream, a refresh trace, a model-visible projection, and an evaluator-only Oracle. Source payloads record authorization, participants, modality, spatial scope, availability, event-time basis, and validity. The validator checks source/evidence field equality, digest closure, dependency acyclicity, an eligibility check on triggers, and strict timestamps with explicit JST offsets. This is synthetic-data provenance under a declared authoring process; hashes establish artifact integrity, not that a source was honestly labeled or that real-world sensing occurred. Figure~\ref{fig:evidence-plane-overview} summarizes the separation between observer-visible context, evaluator-only labels, and the audit boundary.

\begin{figure}[t]
\centering
\includegraphics[width=\linewidth]{figures/part1_evidence_plane_overview.pdf}
\caption{Evidence planes and audit boundary. Source and refresh artifacts feed the observer-visible projection; the evaluator-only Oracle remains outside that context, while the gate audit checks reachability, provenance and delivery integrity.}
\label{fig:evidence-plane-overview}
\end{figure}

The spatial substrate is reused from candidate space packages maintained by the deployed project, including topology and geo references. For the sealed corpus, factual subjects, observer identities, evidence claims, queries, and Oracle labels are generated for this experiment; they are not production character conversations or observed user events. The V1.5 coverage layer is separate: its subject identifiers follow the deployed project canon, while released locators must be opaque and resolve to real coordinates only through a non-released module.

\subsection{Projection and refresh}
We cross three refresh strategies: \texttt{on\_query}, \texttt{fixed\_interval}, and \texttt{event\_trigger}. We cross these with an evidence-scoped policy and a timestamp baseline with a 1,200-second freshness window (a design constant of the protocol, not a field of any artifact; it is labelled as such here because an earlier revision of the Threats section claimed every number in this draft was artifact-bound or marked derived). The projector consumes observer evidence, source-attested validity, refresh trace, and policy only; it has no world-timeline or Oracle argument. Point support is current only at its observation instant under the evidence-scoped policy. Plans remain plans and do not imply execution. Unknown is represented as absence of admissible support and is labeled only in the evaluator Oracle.

\subsection{Independent unit and the fault-targeted invariance audit}
The intended independent unit is a concrete timeline family. Checkpoints, refresh strategies, projection policies, repeats, and paired perturbations are nested observations. The v3 audit moves an inaccessible state end across the query instant, but \texttt{project\_query} accepts the hidden plane through \texttt{**\_hidden\_world} and discards it. The evaluated projection, context, and request payload thus use identical observer inputs. This checks repeatability on that path, not whether the gate rejects an implementation with the historical leak. A separately executed reference projector, written to reproduce the historical leak semantics, changes in 216 nested arms; no faulty implementation was substituted into the evaluated chain. The visible-signal control checks input sensitivity only.

Withholding inaccessible data is a valid architectural boundary. The adequacy gap is the absence of an end-to-end fault injection that makes this same gate fail. Future hidden-plane perturbation checks must change the hidden world artifact before the real observer extraction boundary while preserving observer observations. Source-attested validity mutations instead change observable evidence and require an appropriate output reaction. Writing hidden expiry into declared \texttt{valid\_until} changes the information policy and is not a hidden-world noninterference test. Input differences and expectations must be reported per plane; zero observer-input differences alone do not imply vacuity: the v3 equality was vacuous because no faulty implementation was present in the evaluated chain for the gate to reject. Separating a mutant that never reaches the assertion from one the assertion fails to reject is the stillborn-versus-survived distinction used in benchmark mutation auditing (Bhadra, 2026), and is applied here rather than proposed.

\subsection{Allocation}
The development batch contained two timeline families in each of three classes: \path{home_acoustic_scope_exit}, \path{home_visual_occlusion}, and \path{school_schedule_vs_presence}. The subsequent allocation froze twelve family identifiers, three in each of four classes: \path{public_private_access}, \path{coverage_edge_vs_unobserved_inside}, \path{cross_day_history_retention}, and \path{recipient_scoped_communication}. The allocation commitment is \texttt{\seqsplit{75186e728faac41f73660eb78e20d2515bbf10063cb1103f1aa995b61db41631}}. We apply SHA-256 to the bytes from \texttt{rfc8785.dumps} (version 0.1.4) of the parsed \path{sealed_allocation_v1.json} object with only \texttt{allocation\_sha256} removed. The reproduced canonical pre-image is supplied as \path{peer_review_revision_v1/allocation_preimage.jcs.json}. The pre-image retains the historical analysis-unit wording for byte-level provenance; the present manuscript uses the more precise family-level description. Shared generation precludes estimating variation across independent authors.

\section{Results}
All six development families passed the no-model contract gate, yielding 144 conditions. The frozen stage crossed 12 families, 3 refresh strategies, 2 policies, and 4 checkpoints for 288 conditions; this total is the \texttt{projection\_conditions} field in \path{sealed_timeline_v1/validation_report.json}. Frozen validation reported zero errors and zero answer-model calls after code-derived sampler provenance replaced a placeholder. The v3 report records 216 reference-sensitive nested arms across nine families: 9 distinct hidden-window perturbations, one per family, each repeated across 24 nested arms under one mutation operator and one fault class. The evaluated projection/context/payload digests are identical. Because the hidden parameter was discarded, we withdraw its noninterference PASS interpretation. Nine families supplied visible-signal controls, each attributed to one \texttt{evidence\_id}. The separate recipient gate recorded 27 isolation checks across three communication families; these scoped checks do not validate the hidden-state gate or answer quality.

A three-family validity diagnostic re-digested and revalidated attested variants in all 72 conditions; this total is the \texttt{projection\_conditions} field in \path{sealed_validity_v2/validation_report.json}. Its artifact reports 62 column-changing arms and 50 arms for which the gate's expectation predicate required a reaction. An arm-level re-derivation that reproduces all five frozen aggregates exactly attributes the remaining 12 to the re-attestation clamp: the variant builder floors each moved expiry at one minute after its own observation, so an expiry can still land before the query instant inside a band where the expectation predicate deliberately requires nothing. All 12 fall in that band and none is unexplained, so the gap reflects a predicate narrower than the reaction surface rather than a set of erroneous changes. The hidden side records 24 comparisons, of which only 10 are informative: the window filter admits rows irrespective of snapshot membership while the reference projector skips unadmitted rows, so in the other 14 the perturbed row never reaches the query snapshot and neither side can respond. Reference sensitivity covers all 10 informative comparisons, but that figure is an identity over this sample and not a rate: across all 24 rows the flag that makes a comparison informative and the flag that makes it reference-sensitive are perfectly collinear, so no input could have yielded 9 of 10. Sensitivity is confined to the single family that carries a hidden window at all; the honest power denominator is 10, not 24, and the fault class is exercised in one of three families. The two sides are not structurally matched---attested variants change consumed evidence bytes, whereas hidden variants change discarded arguments---so they are not two arms of one comparison; admitting a comparison before treating its arms as evidence is the discipline reported by Fu et al. (2026). This shows declared expiry can affect projection, but establishes neither hidden-state noninterference nor superiority over inferred invalidation.

\begin{sloppypar}
\paragraph{V1.5 coverage layer.} The adapter produced twelve frozen SceneFrame lineages from nine admitted project packages spanning homes, schools/institutions, a public indoor venue, and an outdoor commercial corridor, modeled on multiple Tokyo locations. Candidates required explicit source, observer scope, validity, and Oracle authoring. The twelve lineages covered 216 projection conditions (12 lineages x 18 nested arms each; a different quantity from the 216 reference-sensitive arms in the sealed-corpus audit) with zero model calls. Historical reports recorded PASS for contract, package-scope, hidden-plane perturbation, validity, and visible controls in the lineages that produced candidates; the thirteenth zero-candidate package is INCONCLUSIVE. Those hidden-plane PASS labels do not overcome the discarded-hidden-input limitation and are not used as noninterference evidence. Original frames remain separate from the sealed corpus; sunshine-city and Arakawa fixtures remain development-only, and a zero-candidate package records an adapter gap. This documents coverage and protocol transfer, not population-level geographic generalization.
\end{sloppypar}

\subsection{External replay and the withdrawn fixture inference}
At GateMem commit \texttt{\seqsplit{603f9f4b4ba4b77f043c20f85687fa016fd720b0}}, the household TF--IDF policy baseline with stub providers ran without model calls. The supplying artifact is \path{results/gatemem_rag_policy_stub_household/summary.json}, whose \texttt{n\_checkpoints} is 552, whose \texttt{n\_privacy}, \texttt{n\_safety} and \texttt{n\_utility} are 184 each, and whose four leakage-rate fields supply every rate quoted below; an earlier revision of this draft quoted those numbers without naming the file they come from. The stub replay contains 552 checkpoint records. Within the labeled subsets, context leakage was reported for 126/184 privacy checkpoints (68.48\%) and 141/184 deletion-labeled checkpoints (76.63\%). Both numerators are derived---each is the reported rate multiplied by its denominator of 184---because the artifact records the rates and not the counts. Final-answer leakage was 0\%. The replay report does not reconcile the remaining 184 records or identify episode- or harness-level clustering, so these are checkpoint-level counts within one external case study rather than cross-system evidence. This reflects a stub configuration, not learned-model safety.

Our paired PIN fixture changed episode text from \texttt{2864} to \texttt{9999} but retained \texttt{2864} in both \texttt{leak\_targets} and \texttt{judge\_spec.not\_include}. The new value remained in the prompt, while the old-target scorer returned 0.0 instead of 1.0. Re-scoring that identical prediction with both target fields updated to \texttt{9999} restored detection to 1.0. Thus our transformation introduced a target/data mismatch; it does not demonstrate an upstream oracle defect on a valid renamed fixture. The visible-query control still reports leakage. This is one external case study, not a cross-system estimate.

\subsection{What the control's two re-issues found}
The control has since been re-issued twice as a self-testing gate, and each re-issue found a conjunct in its predecessor that could not fail. The first found that the earlier version's headline conjunct compared a context digest against itself, because both cases load the same prediction. The second found a further conjunct half that was likewise a theorem of the script's own control flow: the stale branch is an unmutated deep copy, so its field diff against the frozen record is empty by construction and no input could make it non-empty---even though this manuscript had by then begun to claim that no conjunct survived as a theorem. Both are now recorded as disclosures rather than as checks, and the paired form that replaced the second one diffs the two scored checkpoints against each other, which would catch reference aliasing between the branches. The current version carries thirteen checks and fifteen perturbations of its own inputs, each of which flips the verdict, with no vacuous perturbation remaining. The full narration, including the revision of this draft that asserted completeness between the two re-issues and had to be re-issued to find its own error, is in Appendix~\ref{app:reissue-history}: it belongs there because it is a history of this audit rather than a result about the system being audited.

\subsection{The matcher's word-boundary guard}
Static probes of the upstream matcher, recorded with the scorer's digest---the digested file is the upstream scorer at the pinned commit, and its name and full digest are in the gate-status register named under Data and Code Availability rather than here---show it is case-insensitive, whitespace-normalising, and free of Unicode compatibility folding. An earlier version inferred from those probes that a co-transformed representation change would leave the verdict unchanged, and that representation-equivalence operators would therefore yield no valid positive; that inference is withdrawn. The nine probes were single-point assertions, and two of them already held the refuting facts side by side without ever comparing them, so the block had no detection power over the property it was cited for. The replacement evaluates seven paired relations, each on an ASCII arm and a co-transformed arm, and asserts the relation between the two verdicts rather than either verdict alone. Three of the six certified equivalence operators break that relation, in both directions.

The matcher wraps a target in word-boundary guards only when the normalised target fullmatches a character class that the script extracts from the upstream source and records in the artifact; full-width digits and the comma lie outside that class, so a co-transformed target silently loses the guard and over-detects---\texttt{9999} does not hit inside \texttt{99990}, while its full-width twin does hit inside the full-width counterpart---and an unchanged ASCII target under-detects in a CJK-adjacent context, because CJK letters count as word characters. Case and whitespace, the ideographic space included, are the only representation changes certified benign for this scorer. The consequence for the next phase is the reverse of the one previously recorded: a benign-transformation class cannot be assumed, it must be certified per operator against the scorer that consumes it. This remains a static, source-grounded observation about a string matcher. It is not a mutation kill and contributes nothing to any kill-rate denominator, and it is not gate detection power. It is also not, by itself, evidence that any shipped fixture is mislabelled; the paragraph below reports what the same mechanism does on the corpus the benchmark publishes, and deliberately does not adjudicate it.

\paragraph{The same mechanism measured on the published corpus.} Because the mechanism above was established on
strings we wrote ourselves, we re-measured it on the corpus the benchmark ships, model-free and read-only, at the
same pinned commit. Across 91 episodes and 2218 checkpoints
there are 7652 target instances in three fields (\texttt{leak\_targets},
\texttt{judge\_spec.include}, \texttt{judge\_spec.not\_include}). In 59 of them the
target occurs verbatim in the episode text visible as of that checkpoint, after whitespace and case folding, and the
official matcher nonetheless returns no hit. 14 of those are pure-digit targets of
three characters or fewer found inside a longer number---\texttt{26} inside a year, \texttt{94} inside a longer
identifier---and for those the word-boundary guard is behaving correctly, so counting them as misses would overstate
the finding; they are excluded and reported separately. The remaining 45 are
identifier prefix or suffix probes (44 prefixes and 1 suffix), of which 34 sit in
\texttt{leak\_targets} (24 safety and 10 privacy
checkpoints, 1.31 percent of that field's 2601 instances; the 24, the 10 and the 1.31 percent are reductions of the named corpus fields rather than fields in their own right, and are marked as derived), 10 in \texttt{judge\_spec.not\_include} and 1
in \texttt{judge\_spec.include}. They have the form \texttt{TC-20}, \texttt{JRX-77}, \texttt{ORX-80}, \texttt{LPX-91}, each a proper prefix of a longer identifier
that the episode text does contain. The upstream documentation defines \texttt{leak\_targets} as ``protected
strings or patterns for leakage evaluation'' and nowhere states a whole-token requirement, while the guard that
produces these misses appears in no documentation we found.

We do not adjudicate this. If \texttt{leak\_targets} are meant to detect the presence of a protected string in
model-visible text, then a literal that is present and reported absent is a false negative in the official leak
metric. If they are meant to be whole-token matches, the guard is correct and the fixture chose a target that is a
prefix of a longer identifier, which is a labelling matter. Deciding between the two would require a classification
rule, and any rule we wrote now would be written after seeing the data, so it could not be pre-registered; both
readings are therefore recorded per instance in the artifact and neither is asserted. Four further limits apply.
The first is that an instance count is not a consequence. Upstream scores a checkpoint by disjoining
\texttt{any\_pattern\_hits} over a \emph{set} of patterns---\texttt{leak\_targets} for safety, and the union of
\texttt{not\_include} with \texttt{leak\_targets} for privacy---so a suppressed target changes nothing when
another pattern in the same set still hits. Recomputed per checkpoint, 58 of the 59 suppressed instances fall in
a set that is actually scored; the 59th sits in \texttt{judge\_spec.include} on a privacy checkpoint, a field the
scorer never reads for that query type, and it is reported separately rather than dropped, so that all 59 are
accounted for. Those 58 span 40 checkpoints, and under the reading on which a verbatim-present target should have
hit, the number of those checkpoints whose context-side verdict would change is \textbf{zero}: in every one,
another target in the same set still hits. The end-to-end verdict disjoins the answer side with the context side
and this measurement has no model outputs, so it bounds the context-side component and makes no claim about the
disjunction. The script that computes this carries a self-test in which one fixture \emph{must} flip, because a
zero reported by a detector incapable of producing a non-zero is not a result. The counts above are therefore a
statement about the matcher and the zero is the statement about scores; reporting the first without the second
would invite exactly the inference this paper elsewhere declines to license.
No model output is scored: the leak metric runs on a retrieved prompt context, whereas this measurement runs on the
source episode text, so what is established is that the matcher would not see these literals in this form, not that
any published score is wrong. The result is not a mutation kill and satisfies no stop condition in this paper. And
a plausible-looking companion claim was measured and rejected: of 1006 comma-bearing
targets, 971 give different verdicts from their separator-stripped form, but in
0 of those do both forms occur in the text, so the divergence is the expected
behaviour of two different literals rather than a representation artefact, and it is not reported as a finding.
 The stale-versus-aligned verdict pair still cannot distinguish a fragile oracle from a correct one, because that aligned verdict is entailed by literal substring semantics; the representation-equivalence block can distinguish them, and did. Upstream MIT code and data marked CC-BY-4.0 are not redistributed here (Ren et al., 2026).

\begin{table}[h]
\centering
\caption{No-model authoring results.}
\label{tab:authoring}
\begin{tabular}{lrrr}
\toprule
Stage & Families & Conditions & Model calls\\
\midrule
Development batch & 6 & 144 & 0\\
Frozen authoring & 12 & 288 & 0\\
Validity corpus & 3 & 72 & 0\\
\bottomrule
\end{tabular}
\end{table}

\paragraph{The nine mechanisms.} The identifiers in the first column are what this draft actually contributes,
and they classify assertions rather than mutants. \textbf{M1}: the perturbed input never reaches the artifact
under comparison, because it is discarded upstream of the thing being compared. \textbf{M2}: the asserted
equality is a theorem of the callee's signature, since the parameter carrying the perturbation is swallowed and
never read. \textbf{M3}: the pass threshold is ``non-zero'' rather than a declared coverage floor, so one
sensitive arm passes for the whole corpus. \textbf{M4}: the denominator counts comparisons that are vacuous for
both arms, inflating an apparent detection rate. \textbf{M5}: a conjunct compares a value against itself,
because both operands are loaded from the same source. \textbf{M6}: one half of a conjunct is a theorem of the
auditing script's own control flow, so no input to the audited system can make it false. \textbf{M7}: an empty
input set satisfies the predicate, so a family with no rows passes. \textbf{M8}: the probe tests for the
presence of a string rather than the direction of a claim, so it passes identically before and after the claim
is withdrawn. \textbf{M9}: a battery of single-point probes cannot falsify a relational claim, because no probe
compares one arm against its transformed twin. Published mutation audits register verdict classes for
mutants---killed, survived, stillborn, equivalent, unsupported---which is an axis about the mutant. A reviewed published oracle-side property is expected-value anchoring, which names and measures one mechanism by which an oracle cannot fail (Canedo, 2026); Hill's verifier-side audit separately enumerates five fail-open and vacuous-pass classes. Contract audits of tool-using agent environments check advertised tool surfaces against their implementations, tracing each score's provenance to a verdict (Bellibatlu et al., 2026).
M1 to M9 are an axis about the assertion. Neither enumerates the structurally distinct reasons why a
deterministic, model-free assertion in an evaluation gate could not fail, and neither addresses the recursive
case below; a gate can be clean on the mutant axis and vacuous on the assertion axis. Anchoring is the same
family as M5, whose two operands come from one source, and a near relative of M6, which concerns the auditing
script's own control flow rather than the oracle under test; so what is claimed here is not the first sighting
of that mechanism but its co-occurrence with eight others inside one deterministic gate suite, each entry
bound to a retained artifact and a script digest.

One entry is recursive, and it is the reason this register exists in its present form. The first re-issue of
the target-alignment gate reported that no vacuous conjunct remained. That report was itself an instance of M6
one level up: the claim was carried by a conjunct half that no perturbation reached, so no input could have
refuted it, and a second re-issue was needed to find it. Two re-issues exist and no more, so the ordinals here
are checkable against the two script versions on disk rather than against a reader's patience; an earlier
revision of this paragraph said ``third'' in one place and ``second'' in another, and the compiled PDF carried
both. An instrument that reports the absence of vacuity is
subject to the same failure modes it reports, and the only defence we found was to perturb the completeness
claim itself. The same observation, reached on a different instrument, is published independently (Hill, 2026, \S7.4), so this entry records a convergence rather than a discovery.

\subsection{Exploratory response usability}
After the then-reported context and recipient checks, we ran the exploratory batch with temperature 0 and a 1,024-token output limit. All 288 calls completed with zero retries; that zero is a property of the run configuration and is labelled a design constant here rather than presented as an artifact field. The primary status is mutually exclusive: 217 \emph{usable}, 68 \emph{empty}, and 3 \emph{nonempty truncated} (artifact fields \texttt{usable}, \texttt{empty}, and \texttt{nonempty\_truncated}). Separately, 71 calls ended with \texttt{finish\_reason=length} (artifact field \texttt{truncated}; 24.65\%). All 68 empty responses are also length-terminated, so the 68 empty and 3 nonempty-truncated arms are subsets of the 71 length-terminated calls; their union is exactly 71. These columns must not be added or interpreted as model-quality scores.

\begin{table}[h]
\centering
\caption{Descriptive response usability by batch, refresh strategy, and projection policy.}
\label{tab:usability}
\small
\begin{tabular}{lrrrrr}
\toprule
Group & Nested arms & Usable & Empty & Length finish & Usable rate \\
\midrule
Overall & 288 & 217 & 68 & 71 & 75.35\% \\
\texttt{event\_trigger} & 96 & 75 & 19 & 21 & 78.12\% \\
\texttt{fixed\_interval} & 96 & 74 & 22 & 22 & 77.08\% \\
\texttt{on\_query} & 96 & 68 & 27 & 28 & 70.83\% \\
\texttt{evidence\_scoped} & 144 & 117 & 27 & 27 & 81.25\% \\
\texttt{timestamp\_ttl} & 144 & 100 & 41 & 44 & 69.44\% \\
\bottomrule
\end{tabular}
\end{table}

The family-level view is in Appendix Table~\ref{tab:family-usability}. It is kept there because this draft reports no inference over it: the rows are nested arms, the intended independent unit is the family, and arm-level rates would be pseudoreplicated. It is descriptive only: each row contains 24 nested arms, while the intended independent unit remains the family. We report no confidence intervals or significance tests at the arm level: arms are nested within families, so arm-level inference would be pseudoreplicated (Hurlbert, 1984). All rates in Tables~\ref{tab:usability} and~\ref{tab:family-usability} are point descriptions; any future inference must use the family as its unit. The communication families have the highest observed usability in this batch, whereas several access and coverage families have more empty responses; no strategy, policy, or family contrast is interpreted as causal before semantic rating.



\section{Discussion and Limitations}
The results document executable evidence representation and deterministic replay on frozen package clusters. The self-audit withdraws an inadequate gate interpretation and an apparent external oracle defect confounded by stale target labels. The exploratory batch does not estimate temporal awareness, evidence compliance, target accuracy, or mechanism effects; semantic ratings remain deliberately pending. Families share one authoring procedure and package clusters, and nested observations do not increase the family-level sample size. No strategy or policy contrast is treated as causal. A zero-evidence or opaque-entity control is required before new model calls to separate evidence use from prior knowledge of recognizable venues.

\paragraph{What the audit invalidated.}
The two v2 checks, the v3 evaluated hidden-input check, and the twelve V1.5 branch hidden-plane reports of Section 4 all share the discarded-hidden-input defect; none of their PASS labels is cited as noninterference evidence. Their artifacts remain for audit continuity. The 288-call manifest also records a superseded script whose present hash differs from its frozen hash: \path{run_context_twin_gate_v2.py} changed from 4,184 to 4,485 bytes (\texttt{04a3f0b6\ldots9fdd74} to \texttt{e4dfb9f8\ldots75a69}); the drift is confined to the module docstring, which now carries the withdrawal notice, and no executable line changed. The new audit records both complete hashes as DRIFT. Historical manifests are preserved rather than rewritten, so a third-party integrity check of the paid batch will still report DRIFT; descriptive response counts do not repair that provenance gap. An earlier internal authorization-leak finding was also overstated, as corrected below. Table~\ref{tab:gate-status} lists every gate cited in this draft together with its status and the defect that produced it.

\begin{table}[h]
\centering
\caption{Gate status and known-bad register. VOID means the recorded verdict must not be cited; SUPERSEDED means a later script fixes a defect that made the earlier verdict uninformative; REF-ONLY PASS is the machine status \texttt{PASS\_REFERENCE\_ONLY} and is used only for rows whose artifact actually records it. A row that instead reads ``PASS as recorded'' names an artifact carrying no such flag, and the reclassification beside it is this draft's reading rather than a machine status. Full script digests are in the gate-status register named under Data and Code Availability.}
\label{tab:gate-status}
\scriptsize
\setlength{\tabcolsep}{3pt}
\begin{tabular}{@{}p{0.055\textwidth}p{0.265\textwidth}p{0.15\textwidth}p{0.42\textwidth}@{}}
\toprule
Id & Gate / artifact & Status & Defect or limitation \\
\midrule
M1 & \path{run_twin_gate_sealed_v2.py} & VOID & perturbed a world object no function received \\
M1 & \path{run_context_twin_gate_v2.py} & VOID + DRIFT & same defect; recorded digest no longer matches (docstring-only change) \\
M1, M2 & \path{twin_v4/twin_gate_report_v4.json} & REF-ONLY PASS (withdrawn) & hidden plane swallowed at \path{validate_observer_v2.py:642}; equality is a theorem of the callee signature; all sensitivity comes from a bypass reference projector; no fault was injected into the evaluated chain. \texttt{PASS\_REFERENCE\_ONLY} is v4's machine status; the superseded \path{twin_v3/twin_gate_report_v3.json} recorded an unqualified \texttt{PASS} for the same defect, and an earlier revision of this row cited v3 while quoting v4's status \\
M1, M2 & 12 x \path{v1_5_admitted_*/twin_v3/} & PASS as recorded; reclassified here as reference-only (withdrawn) & same discarded-hidden-input path. These twelve artifacts record an unqualified \texttt{PASS} and carry no reference-only flag, so unlike the \path{twin_v4} row above the reclassification is this draft's reading of them rather than a machine status; that reading is carried by a sidecar, \path{twin_v3/INTERPRETATION_WITHDRAWN_2026-10-05.md}, which names these twelve by pattern, and the bytes are not rewritten. The sidecar's own stated delivery mechanism is that it is read alongside its directory, which holds for the directory containing it and not for these twelve, whose directories hold only their two JSON files: a reader who opens one of the twelve sees an unqualified \texttt{PASS} with no warning beside it. A warning whose delivery mechanism does not reach what it warns about is the failure mode this paper audits, so the gap is stated here rather than closed silently, and an earlier wording placed the sidecar beside bytes it does not sit beside. The 13th zero-candidate package is INCONCLUSIVE \\
M3, M4 & \path{validity_v1/validity_gate_report_v1.json} & VALID, LIMITED & the attested side does run the evaluated chain; but the hidden denominator is inflated (24 comparisons, 10 informative), the PASS threshold is merely non-zero with no coverage floor, and there is no INCONCLUSIVE branch \\
M5 & \path{audit_gatemem_target_alignment.py} (v1) & SUPERSEDED & its headline conjunct compared a value against itself: both cases load the same prediction, so the equality could not fail \\
M6, M9 & \path{audit_gatemem_target_alignment_v2.py} & SUPERSEDED & 11 of 11 self-test mutations flipped, but a second conjunct half was still a theorem of the script's own control flow, and its recorded consequence asserted a representation-equivalence stability that v3 refutes \\
-- & \path{audit_gatemem_target_alignment_v3.py} & REPRODUCED & 15 of 15 self-test mutations flip the verdict; both former theorems demoted to disclosures; 3 of 6 certified equivalence operators break verdict preservation, in both directions \\
M7 & \path{audit_confirmatory_query_readiness.py} (v1) & SUPERSEDED & a family with no query rows passed with \texttt{query\_count: 0}, and the leakage count duplicated the placeholder count \\
-- & \path{audit_confirmatory_query_readiness_v2.py} & FAIL & all 84 confirmatory queries are still placeholders; detection power PARTIAL with four unexercised branches. This is the state a measurement phase would start from, and it is not a startable state \\
M8 & \path{run_paper_claim_consistency_audit.py} (v1) & NO DET.\ POWER & passed identically on the pre- and post-withdrawal manuscripts; its checks test for the presence of strings, not the direction of a claim \\
\bottomrule
\end{tabular}
\end{table}

\paragraph{Boundaries of the contract claim.}
The 288-condition corpus behind the answer batch uses point support throughout, so the source-attested validity branch of the projector---the only branch that can carry an expiration---was exercised by none of those conditions, and under the point rule the evidence-scoped policy reduces to ``current only at the observation instant''. A separate three-family corpus now exercises that branch and passes strict validation in all 72 of its conditions, but it has no model calls attached to it yet, so no model-level validity result exists. The projection contract also emits no unknown column: the schema defines only \texttt{current} and \texttt{history}, and across all 288 rebuilt prompt contexts the \texttt{unknown} field is empty, even though the answer prompt tells the model that such a column exists and the Oracle labels unknown checkpoints with reason codes in all twelve families. Unknown calibration is therefore measured under a stronger condition than a supplied unknown list, and the contract must be described as evaluator-side rather than projected. On authorization we correct an earlier internal finding: the contract validator already rejects the two mis-addressed artifact kinds injected here (an unauthorised source reference and a non-recipient observer). Under those two probes, none of the 12 guarded families (\texttt{guarded\_families\_tested} in \path{twin_v4/twin_gate_report_v4.json}; each of the twelve per-family artifacts records 0 and only the root report carries the aggregate) reached context without authorization; we therefore report no leak for the probed kinds only. A projector guard has also been added: it is a no-op on the frozen corpus (aggregate column-assignment digest \texttt{e96e2b81\ldots604d8e}, recorded in \path{FINAL_PAPER_PLAN_2026-10-03.md} \S13.2) but can fail—re-addressing one communication evidence row to a non-recipient drops its rendered row count from 1 to 0. Spatial scope for direct observations remains the package validator's responsibility, is outside these probes, and is not covered by any gate reported here; we make no claim about mis-addressing kinds we did not probe.

\paragraph{Missing comparisons.}
No existing memory system is used as a baseline, so we cannot state that source-attested expiry outperforms the inferred invalidation used by deployed temporally-aware memory layers. The only contrast currently in the design is our own freshness policy against our own evidence-scoped policy. The response batch used one model, \texttt{deepseek-flash}, at temperature 0 with a 1,024-token output limit; the artifact records no finer model revision identifier. There is one prompt assembly and no third-party reproduction, so these descriptive results are not treated as transferable. A counterfactual baseline that simply places the full visible context in the prompt is defined in the protocol but not yet evaluated.

\section{Anatomy of vacuous assertions}
\label{sec:anatomy}
The register table above lists nine mechanisms by identifier, beside the gate that exhibited each and the artifact that retains it. The class they belong to is not named here for the first time: an audit of an evidence-bundle verifier publishes a named definition of vacuous pass, maps that definition to CWE entries, and exhibits working forgeries against its own hardened verifier (Hill, 2026). What this section adds is a per-mechanism treatment underneath that published class. This section is the part meant to travel. For every mechanism it gives four things: a structural signature that never mentions this system, the reason an existing discipline does not catch it, a question a reader can put to their own gate without running any of ours, and the smallest example we found that fails. None of it depends on the host project. The stronger claim an earlier wording made here---that none of it depends on the external benchmark or on any count reported elsewhere in this draft---is falsified by four of this section's own minimal examples, which quote the 72-condition validity corpus, the 10 informative comparisons out of 24, the \texttt{2864} fixture value and the nine visible positive controls, and one of which uses the external benchmark's matcher. What travels is therefore the four-element treatment of each mechanism, not the independence of every illustration. Nor are those examples minimal in size: four are transcribed from this package's real artifacts at their original scale, so ``the smallest example we found that fails'' describes the role an example plays in the argument and is not a proof that no smaller one exists.

The four elements are not decoration and they are not interchangeable. The signature says what to look for. The blindness explanation says why a discipline that already has a name for the symptom still misses it, which is the only reason to give a mechanism its own entry rather than folding it into an existing one. The question is the transferable artifact, because a reader who cannot run our gates can still run it. The minimal example is what separates a mechanism from a mood.

\paragraph{M1: the perturbed input never reaches the compared artifact.}
\emph{Signature.} A stage between the injection point and the comparison discards, ignores or overwrites the perturbed plane, so both arms of the comparison are built from the same bytes.
\emph{Why it is missed.} Mutation testing already has a word for this---a mutant no assertion ever sees is \emph{stillborn}---but stillborn mutants are removed from the kill-rate denominator as a bookkeeping step, not investigated as a fault in the harness. The discipline classifies the symptom and then discards it.
\emph{Question.} For every perturbation your suite injects, dump the bytes on both sides of the comparison and diff them. An empty diff means the perturbation is not in the artifact under test, whatever the verdict says.
\emph{Minimal example.} Assert \texttt{normalise(a) == normalise(b)} where the injected difference lies in the part of \texttt{a} that \texttt{normalise} strips.

\paragraph{M2: the asserted equality is a theorem of the callee's signature.}
\emph{Signature.} The property under test is entailed by a type, a schema or a function contract, so no runtime input can violate it.
\emph{Why it is missed.} A check that always passes looks identical, in a pass/fail report, to a check that passes because the system is correct. Test-count and coverage metrics both count it as an executed assertion.
\emph{Question.} For each equality you assert, name the input that would falsify it. If you cannot construct one without changing the callee's signature, you have proved a theorem, not tested a system.
\emph{Minimal example.} Assert that a function annotated to return a non-empty list returns a non-empty list.

\paragraph{M3: the pass threshold is ``non-zero'' rather than a declared floor.}
\emph{Signature.} The predicate is \texttt{count > 0} where the claim being made needs \texttt{count >= k} for some \texttt{k} fixed before the run.
\emph{Why it is missed.} Adequacy criteria in the testing literature are usually stated as coverage obligations, so a gate that reports any detection at all reads as satisfying the criterion. The gap between one and the required number is invisible unless the required number is written down.
\emph{Question.} Find every threshold in your gates and ask whether it was declared before the run or read off the result. Any threshold you cannot date is a threshold you chose afterwards.
\emph{Minimal example.} A sensitivity gate that passes on one responsive arm out of seventy-two, reported as ``sensitivity demonstrated''.

\paragraph{M4: the denominator counts comparisons that are vacuous for both arms.}
\emph{Signature.} A rate is reported over a set of comparisons, some of which cannot discriminate because neither arm could ever produce a difference.
\emph{Why it is missed.} Denominator hygiene is discussed for kill rates, where equivalent mutants are excluded. It is rarely applied to comparison counts, and an inflated denominator makes a gate look more thorough rather than less.
\emph{Question.} For each comparison in a denominator, ask what observation would make its two arms differ. If the answer is ``nothing'', it is not a comparison.
\emph{Minimal example.} Reporting 10 of 24 hidden comparisons as sensitive, when in 14 of the 24 the perturbed row never entered the queried snapshot and so neither arm could react.

\paragraph{M5: a conjunct compares a value against itself.}
\emph{Signature.} Two operands in an equality are loaded from the same source, often through different variable names, so the comparison is reflexive.
\emph{Why it is missed.} Static analysis flags self-comparison in a single expression; it does not usually track that two differently named locals alias one loaded object. Reviewers read the assertion, not the data flow behind it.
\emph{Question.} For every equality, trace both operands to their load site. Two paths from one load is one operand.
\emph{Minimal example.} \texttt{assert case1.context == case2.context} where both cases were built from the same prediction object and only the checkpoint differs.

\paragraph{M6: one half of a conjunct is a theorem of the auditing script's own control flow.}
\emph{Signature.} A conjunct is guaranteed by how the checker was written---a copy that is never mutated, a branch that is never taken, a default that is always applied---rather than by anything about the system under test.
\emph{Why it is missed.} This is M2 one level up: the theorem lives in the auditor, not the audited. A self-test that perturbs the inputs of the system under test cannot reach it, because no input to that system can make the conjunct false.
\emph{Question.} Perturb the checker rather than the checked. If some conjunct survives every perturbation of the audited system, ask what it is actually a function of.
\emph{Minimal example.} Diffing a branch against a record it was deep-copied from, then asserting the diff is empty.

\paragraph{M7: an empty input set satisfies the predicate.}
\emph{Signature.} A universally quantified check over a collection passes when the collection is empty, and emptiness is a failure of collection rather than a property of the system.
\emph{Why it is missed.} \texttt{all()} over an empty iterable is true in most languages, and a gate that reports ``0 failures'' does not distinguish ``nothing failed'' from ``nothing was examined''.
\emph{Question.} For every universal check, run it against an input you know is empty and confirm it fails. A universal that passes on the empty set is not asserting anything.
\emph{Minimal example.} A family with no query rows reporting \texttt{query\_count: 0} and passing a check that no query is malformed.

\paragraph{M8: the probe tests for the presence of a string rather than the direction of a claim.}
\emph{Signature.} A regression probe asserts that some token appears in a document, where the claim at stake is about what the document concludes.
\emph{Why it is missed.} Presence probes are cheap, stable and easy to write, so they proliferate. A token like a figure or a term appears in a retraction and in an assertion of the same magnitude, and the probe cannot tell them apart.
\emph{Question.} For each presence probe, write the sentence that would make it pass while the underlying claim is false. If you can write it easily, the probe is testing vocabulary.
\emph{Minimal example.} Requiring the string \texttt{2864} to appear in a draft that has withdrawn every conclusion drawn from \texttt{2864}.

\paragraph{M9: single-point probes cannot falsify a relational claim.}
\emph{Signature.} The claim is a relation between two arms---equivalence, monotonicity, invariance---but every probe evaluates one arm in isolation.
\emph{Why it is missed.} A battery of single-point probes can be large and every one can pass, and the report reads as thorough. The relation is never evaluated, so nothing in the battery could have contradicted it. This is the shape metamorphic testing exists to address, and a suite that has not adopted it will not notice the absence.
\emph{Question.} For each relational claim, find the probe that evaluates both arms together and asserts something about the pair. If there is none, the claim is untested however many probes pass.
\emph{Minimal example.} Nine static probes each describing one property of a matcher, cited as evidence that a co-transformed input preserves the verdict, where no probe ever compared two verdicts.

\paragraph{What the nine share.} Every one of them is a case where the assertion, not the system, is what makes the check unfalsifiable, and in most of the nine the report a gate produces is identical whether the check has power or not---eight on a per-entry reading, which is marked here as derived rather than quoted from a field, because no artifact supplies it and because its denominator is the one the construct-validity paragraph below declines to fix, a reader who merges the two near-adjacent mechanism pairs getting seven rather than nine. That is why the register binds each entry to a retained artifact and a script digest rather than to a narrative: the only way to tell M1 through M9 apart from a healthy gate is to perturb the gate and watch which checks move. The ninth, the recursive entry, is the observation that this claim about the register is itself an assertion of the kind the register classifies, and Appendix~\ref{app:reissue-history} records what happened when it was not perturbed.

\section*{Threats to Validity}
\paragraph{Construct validity.} The register's nine mechanisms are an axis about assertions, and the claim that they are structurally distinct rests on each having a produced artifact, not on a formal independence proof. Two pairs sit close together---M5 and M6 both concern a conjunct that cannot fail, differing in whether the reason is the compared values or the auditing script's control flow, and M1 and M2 both concern a perturbation that never becomes a difference, differing in whether it is discarded downstream or absorbed by a signature---so a reader who merges them gets seven mechanisms rather than nine, and nothing in this paper decides between those counts. What survives either count is the axis itself. A completeness claim over the nine would require perturbing the completeness claim, which is the recursive entry, and no such perturbation is offered here.

\paragraph{Internal validity.} Every count quoted from an artifact resolves to a field named under Data and Code Availability, and the derivations that matter were recomputed from per-instance data rather than copied from a prior report: the corpus suppression counts, the verdict-flip figure, and the collinearity of the informative-comparison flags. That is a narrower claim than an earlier revision of this paragraph made. The earlier one said that every count in the draft was either read from a named field or marked as derived, and four classes of number falsified it: the replay counts of Section 4, whose supplying artifact this draft did not name anywhere; the per-metric suppression splits, which are reductions of the named corpus fields rather than fields in their own right and were not marked as derived; the 1,200-second freshness window and the zero-retries setting, which are design constants with no supplying field at all; and the boundary-guard digest, quoted in full while the file it digests is named only in the accompanying register. Each is now named, marked as derived, or labelled a design constant at the place it appears, and the register is named for the digest. A universal statement about one's own numbers is exactly the kind of assertion this paper audits, so the narrowed form is stated with its counterexamples rather than replaced silently. Two further classes of number were wrong in earlier revisions and are corrected here rather than quietly dropped: an ordinal that contradicted itself between two paragraphs of the same compiled PDF, and a Chinese sentence whose antecedent made a numerator equal its own denominator. Both are transcription failures of the kind the register's M13 names---M13 being that register's local identifier for a document which copies a value out of another document and then goes stale when the source moves, defined there and not in this draft---which is the reason the accompanying freeze generator now verifies its own output instead of trusting the fields it copied.

\paragraph{External validity.} One system, one external benchmark snapshot at a pinned commit, no model calls in the gates, and a single rater for anything semantic. The mechanisms are candidate transferable classes, not a measured taxonomy: transferring them requires applying the same per-operator certification and the same corpus scan to a second licence-clean scorer, which has not been done. The corpus finding is explicitly not adjudicated between a scorer false negative and a fixture labelling choice, and its verdict-level consequence in this corpus is zero, so it licenses no claim about any published score.

\paragraph{Conclusion validity.} No inferential statistics are reported, so there is no significance claim to threaten. Where a rate could be computed it is reported with its denominator and, where the denominator cannot discriminate, with a statement that the figure is not quotable as a rate. Wilson intervals appear only where a kill-rate denominator would exist, and no such denominator exists yet.

\section{Conclusion and Next Step}
Part I currently supports a model-free contract feasibility check and a reproducible failure audit of one system's own gates. It does not support contract feasibility in the end-to-end sense: the mechanism that would differentiate the contract, source-attested validity, was exercised by none of the model batches, and the noninterference check that would have carried it was vacuous and is withdrawn. Gate mutation testing and reference differentials are established techniques, not a new contribution by themselves. A cross-benchmark protocol of exactly the shape we had planned has already been published (Sun, 2026), so a follow-up measurement study would have to differ from it in subject rather than in method: the object here is one observer-scoped evidence contract and its own gate history, not task-scoped state oracles in general. A prospective measurement study must inject faults into the evaluated execution path, preserve data/target consistency, and examine multiple independently selected public harnesses. A source-attested policy claim separately requires a baseline and conditions under which it can lose. Neither route is established by the current results. Existing corpora and replay tools remain reusable technical artifacts. Any later model study requires new controls, a frozen analysis plan, and an explicitly single-rater, blinded semantic audit with delayed re-rating; no inter-rater reliability is claimed.

\section*{Venue and the scope of this draft}
This draft is written for a reproducibility and evaluation-methodology venue rather than for an empirical software engineering journal, and the reason is a measurement it does not have. Its external object is one benchmark snapshot at one pinned commit, and the per-operator certification and corpus scan of Section 4 have not been repeated on a second licence-clean scorer, so the transferability of the benign-transformation rule rests on one instrument. An empirical journal would reasonably ask how much of the finding is a property of that scorer rather than of scorers, and this draft cannot answer. Repeating the certification on a second scorer is model-free and read-only and is the shortest route to an answer; it was considered and ruled out of scope for this draft, and the venue was chosen with that ruling in hand, because adding a second artifact family to a package whose governance chain was still converging would have cost more confidence than it added. It is therefore a limitation of this draft, not a task this draft owes; the register records the ruling and this sentence previously described it as still pending. What this draft offers instead is a complete audit trail: every count bound to a named artifact field or labelled as derived or as a design constant, every gate bound to a script digest, every superseded artifact retained---in the qualified sense the superseded-artifact index states, namely not rewritten by this tooling rather than recoverable, since this directory holds no version control---and indexed, and every claim this audit withdrew recorded beside the claim that replaced it.

The domain is the reliability of software evaluation infrastructure---whether the instruments that decide if agent-memory and agent-benchmark results mean what their scores say can themselves fail without reporting it. The lineage is the one Section 2 cites: oracle authority and what oracles are judged against (Mughal and Bilal, 2026), fault triggering as distinct from fault detection (Hamidi et al., 2026), mutation testing turned on evaluators and on benchmark rigor (Bhadra, 2026; Yang et al., 2026; Delcourt et al., 2026; Shulepov, 2026), contract audits of agent environments that trace a score to its verdict (Bellibatlu et al., 2026), the oracle-side property that is the same family as one of the nine mechanisms here and a near relative of a second (Canedo, 2026), and the argument that evaluation reporting must carry its methodology and not only its scores (Kapoor et al., 2025).

\section*{Funding}
This research received no specific grant from any funding agency in the public, commercial, or not-for-profit sectors.

\section*{Data and Code Availability}
The draft is backed by the local experiment bundle \texttt{e0\_observation}, and every artifact count quoted in the text resolves to a named field of a named file. The full field-by-field index---which report supersedes which, the digests of the probe and allocation records, and what is deliberately not redistributed---is in Appendix~\ref{app:availability-index}. It is kept out of the body because it is a lookup table, and reading a lookup table linearly is not how anyone uses it.

The public code, opaque synthetic fixtures, aggregate audit records, figure sources and de-identified manuscript derivatives are available at \url{https://github.com/puresky271/observer-evidence-study}. The release tag and commit identifier will be recorded in the accompanying release manifest. Hidden evaluator labels, reversible mappings and the untransformed source manuscripts are not distributed.

Three things belong in the body rather than in an appendix. First, the gate-status register \path{GATE_STATUS_AND_KNOWN_BAD_REGISTER_2026-10-05.md} records the full digest of every gate script and the status of every verdict cited here, and it is the authority for those values rather than this draft; where this draft and the register disagree, the register and the most recent consistency-audit PASS report win. Second, the per-instance corpus file is \emph{not} redistributed: it carries upstream target values and checkpoint identifiers, and is reproducible from the pinned upstream commit by the named script. Third, this draft's own bytes are not releasable yet, which is a condition rather than a detail: a character name from the host project appears four times in each language version and a real place name once more. Publication therefore requires three separately recorded conditions: an opaque-ID transformation applied only to derived copies; a human release decision for the retained real-place-name class; and the scanner result before that decision is recorded (the current derived copies are 0 ADMIT and 2 EXCLUDE). The divergence from the frozen originals must be recorded, and the originals must not be rewritten; until all three conditions are met this file is not releasable, while the accompanying public repository can expose the tooling, schemas, synthetic opaque fixtures and figure scripts.

\section*{References}
\paragraph{Writing assistance disclosure.} Local claim reconciliation used the project writing and critical-thinking workflows; this assistance is not scientific evidence or human verification.
\sloppy
\begin{enumerate}[leftmargin=*]
\item Wu, Di, et al. ``LongMemEval: Benchmarking Chat Assistants on Long-Term Interactive Memory.'' ICLR (2025). arXiv:2410.10813. \url{https://arxiv.org/abs/2410.10813}
\item Ahn, Jaewoo, Taehyun Lee, Junyoung Lim, Jin-Hwa Kim, Sangdoo Yun, Hwaran Lee, and Gunhee Kim. ``TimeChara: Evaluating Point-in-Time Character Hallucination of Role-Playing Large Language Models.'' Findings of the Association for Computational Linguistics: ACL 2024, pp. 3291--3325 (2024). DOI: 10.18653/v1/2024.findings-acl.197. \url{https://aclanthology.org/2024.findings-acl.197/}
\item Tallam, Krti. ``Partial Evidence Bench: Benchmarking Authorization-Limited Evidence in Agentic Systems.'' arXiv preprint arXiv:2605.05379 (2026). \url{https://arxiv.org/abs/2605.05379}
\item Snodgrass, Richard T., and Ilsoo Ahn. ``A Taxonomy of Time Databases.'' Proceedings of the 1985 ACM SIGMOD International Conference on Management of Data. DOI: 10.1145/318898.318921.
\item TSQL2 Language Specification Committee. ``The TSQL2 Language Specification (Version 1.2).'' SIGMOD Record 23(2) (1994). DOI: 10.1145/181550.181562.
\item Snodgrass, Richard T., ed. ``The TSQL2 Temporal Query Language.'' Kluwer Academic Publishers (1995). DOI: 10.1007/978-1-4615-2289-8.
\item Jensen, Christian S., Richard T. Snodgrass, and Michael D. Soo. ``Extending Existing Dependency Theory to Temporal Databases.'' IEEE Transactions on Knowledge and Data Engineering 8(4) (1996): 563--582. DOI: 10.1109/69.536250.
\item Dyreson, Curtis E., and Richard T. Snodgrass. ``Valid-Time Indeterminacy.'' Proceedings of the Ninth International Conference on Data Engineering (ICDE 1993). DOI: 10.1109/ICDE.1993.344048.
\item Microsoft. ``Temporal table conventions'' (system-versioned tables, \texttt{FOR\_SYSTEM\_TIME AS OF}, half-open validity interval). \url{https://learn.microsoft.com/en-us/sql/relational-databases/tables/temporal/overview}
\item W3C. ``SSN --- Semantic Sensor Network Ontology'' (Recommendation, 19 October 2017): \texttt{phenomenonTime} and \texttt{resultTime}. \url{https://www.w3.org/TR/vocab-ssn/}
\item Rasmussen, Preston, Pavlo Paliychuk, Travis Beauvais, Jack Ryan, and Daniel Chalef. ``Zep: A Temporal Knowledge Graph Architecture for Agent Memory.'' arXiv preprint arXiv:2501.13956 (2025). \url{https://arxiv.org/abs/2501.13956}
\item Shi, Zitong, Yixuan Tang, and Anthony Kum Hoe Tung. ``A-TMA: Decoupling State-Aware Memory Failures in Long-Term Agent Memory.'' arXiv preprint arXiv:2607.01935 (2026). \url{https://arxiv.org/abs/2607.01935}
\item Zhang, Yehang, et al. ``WorldLines: Benchmarking and Modeling Long-Horizon Stateful Embodied Agents.'' arXiv preprint arXiv:2606.18847 (2026). \url{https://arxiv.org/abs/2606.18847}
\item Chao, Hanxiang, Yihan Bai, Rui Sheng, Tianle Li, and Yushi Sun. ``STALE: Can LLM Agents Know When Their Memories Are No Longer Valid?'' arXiv preprint arXiv:2605.06527 (2026). \url{https://arxiv.org/abs/2605.06527}
\item Zhou, Yan, Yue Ouyang, Kaiyang Zheng, and Suncheng Xiang. ``TEPA: Revoking Stale Memories for Conflict-Robust Language Agents.'' arXiv preprint arXiv:2608.07429 (2026). \url{https://arxiv.org/abs/2608.07429}
\item Panthi, Sugam, Muhaiminul Yeamin, Siyan Luo, and Rabab Abdelfattah. ``Memory Consolidation Flattens the Temporal Shape of User Facts.'' arXiv preprint arXiv:2609.36457 (2026). \url{https://arxiv.org/abs/2609.36457}
\item Ouyang, Jie, et al. ``HoH: A Dynamic Benchmark for Evaluating the Impact of Outdated Information on Retrieval-Augmented Generation.'' arXiv preprint arXiv:2503.04800 (2025). \url{https://arxiv.org/abs/2503.04800}
\item Thorne, James, et al. ``FEVER: a large-scale dataset for Fact Extraction and VERification with Unstructured and Structured sources.'' NAACL 2018. DOI: 10.18653/v1/N18-1074. \url{https://arxiv.org/abs/1803.05355}
\item Zhou, Xinyu, Chang Jin, Carsten Eickhoff, Zhijiang Guo, and Seyed Ali Bahrainian. ``When Silence Is Golden: Can LLMs Learn to Abstain in Temporal QA and Beyond?'' arXiv preprint arXiv:2602.04755 (2026). \url{https://arxiv.org/abs/2602.04755}
\item Kirichenko, Polina, Mark Ibrahim, Kamalika Chaudhuri, and Samuel J. Bell. ``AbstentionBench: Reasoning LLMs Fail on Unanswerable Questions.'' arXiv preprint arXiv:2506.09038 (2025). \url{https://arxiv.org/abs/2506.09038}
\item Geifman, Yonatan, and Ran El-Yaniv. ``Selective Classification for Deep Neural Networks.'' NeurIPS 2017. arXiv:1705.08500. \url{https://arxiv.org/abs/1705.08500}
\item Hurlbert, Stuart H. ``Pseudoreplication and the Design of Ecological Field Experiments.'' Ecological Monographs 54(2) (1984): 187--211. DOI: 10.2307/1942661.
\item Goguen, Joseph A., and Juan Meseguer. ``Security Policies and Security Models.'' Proceedings of the IEEE Symposium on Security and Privacy 1982. DOI: 10.1109/SP.1982.10014.
\item Barthe, Gilles, Pablo R. D'Argenio, and Timothy Rezk. ``Secure information flow by self-composition.'' Proceedings of the 17th IEEE Computer Security Foundations Workshop 2004. DOI: 10.1109/CSFW.2004.1310735.
\item Hrițcu, Cătălin, John Hughes, Benjamin C. Pierce, Antal Spector-Zabusky, Dimitrios Vytiniotis, and Arthur Azevedo de Amorim. ``Testing noninterference, quickly.'' ACM SIGPLAN Notices 48(9) (Proceedings of ICFP 2013). DOI: 10.1145/2544174.2500574.
\item Beer, Ilan, Shoham Ben-David, Cindy Eisner, and Yoav Rodeh. ``Efficient Detection of Vacuity in Temporal Model Checking.'' Formal Methods in System Design 18(2) (2001). DOI: 10.1023/A:1008779610539.
\item DeMillo, Richard A., Richard J. Lipton, and Frederick G. Sayward. ``Hints on Test Data Selection: Help for the Practicing Programmer.'' IEEE Computer 11(4) (1978). DOI: 10.1109/C-M.1978.218136.
\item Chen, T. Y., T. H. Tse, and Z. Q. Zhou. ``Fault-based testing without the need of oracles.'' Information and Software Technology 45(1) (2003). DOI: 10.1016/S0950-5849(02)00129-5.
\item Ribeiro, Marco Tulio, Tongshuang Wu, Carlos Guestrin, and Sameer Singh. ``Beyond Accuracy: Behavioral Testing of NLP Models with CheckList.'' ACL 2020. DOI: 10.18653/v1/2020.acl-main.442. \url{https://aclanthology.org/2020.acl-main.442/}
\item Poliak, Adam, Jason Naradowsky, Aparajita Haldar, Rachel Rudinger, and Benjamin Van Durme. ``Hypothesis Only Baselines in Natural Language Inference.'' Proceedings of the Seventh Joint Conference on Lexical and Computational Semantics (SEM 2018). DOI: 10.18653/v1/s18-2023.
\item Garware, Chaitanya Vilas, and Sharif Noor Zisad. ``When the Ruler is Broken: Parsing-Induced Suppression in LLM-Based Security Log Evaluation.'' arXiv preprint arXiv:2605.07293 (2026). \url{https://arxiv.org/abs/2605.07293}
\item Kapoor, Sayash, Benedikt Stroebl, Peter Kirgis, Nitya Nadgir, Zachary S. Siegel, et al. ``Holistic Agent Leaderboard: The Missing Infrastructure for AI Agent Evaluation.'' arXiv preprint arXiv:2510.11977 (2025). \url{https://arxiv.org/abs/2510.11977}
\item Ren, Zhe, Yibo Yang, Yimeng Chen, Zijun Zhao, Benshuo Fu, Zhihao Shu, Bingjie Zhang, Yangyang Xu, Dandan Guo, and Shuicheng Yan. ``GateMem: Benchmarking Memory Governance in Multi-Principal Shared-Memory Agents.'' arXiv preprint arXiv:2606.18829 (2026). DOI: 10.48550/arXiv.2606.18829. \url{https://arxiv.org/abs/2606.18829}
\item Xiong, Wenjun, Shengtao Zhang, Shangding Gu, Bo Tang, Zhiyu Li, Feiyu Xiong, Ying Wen, and Muning Wen. ``TRACE: Governing Memory Validity in Evolving Multi-Agent Systems.'' arXiv preprint arXiv:2609.33517 (2026). \url{https://arxiv.org/abs/2609.33517}
\item Yang, Jiayuxuan, Jie M. Zhang, Yiling Lou, and Zhenpeng Chen. ``Mubric: Mutation Testing-Guided Rubric Generation for LLM Evaluation.'' arXiv preprint arXiv:2609.37322 (2026). \url{https://arxiv.org/abs/2609.37322}
\item Delcourt, Kevin, Meriem Ben Chaaben, Abdelhamid Rouatbi, Luciano Marchezan, and Houari Sahraoui. ``Breaking Models to Test the Judge: A Mutation Testing Approach for Semantic Evaluators of Domain Class Diagrams.'' arXiv preprint arXiv:2608.14315 (2026). \url{https://arxiv.org/abs/2608.14315}
\item Bhadra, Meet. ``GateTruth: Auditing the Rigor of RTL Design Benchmarks via Mutation Testing.'' arXiv preprint arXiv:2608.12635 (2026). \url{https://arxiv.org/abs/2608.12635}
\item Shulepov, Ilya. ``Mutation Testing for Reproducibility Safeguards in Machine Learning Research Software: An Empirical Study.'' arXiv preprint arXiv:2608.27100 (2026). \url{https://arxiv.org/abs/2608.27100}
\item Canedo, Arquimedes. ``Oracles That Cannot Fail: Anchoring and the Expectation That Moves With the Fault.'' arXiv preprint arXiv:2608.17214 (2026). \url{https://arxiv.org/abs/2608.17214}
\item Hill, Erik. ``Four Ways to Forge a Bundle My Own Verifier Calls Clean: Refusal-Site Mutation Testing of an Evidence-Bundle Verifier.'' arXiv preprint arXiv:2608.26183 (2026). \url{https://arxiv.org/abs/2608.26183}
\item Bellibatlu, Rohith Reddy, Zichong Wang, and Wenbin Zhang. ``Do Agent Benchmarks Do What They Say? An Executable-Contract Audit of Tool-Using Agent Environments.'' arXiv preprint arXiv:2609.37315 (2026). \url{https://arxiv.org/abs/2609.37315}
\item Fu, Hao, Baiting Zhu, Minglei Chen, Yinjie Huang, and Shuai Ding. ``Verify, Don't Trust: Agentic Model Development for Video Discovery Retrieval at Scale.'' arXiv preprint arXiv:2609.21257 (2026); ACM reference format in the preprint names KDD '27. \url{https://arxiv.org/abs/2609.21257}
\item Sun, Shengyao. ``Mutation Testing of Task-Scoped State Oracles in Software-Agent Benchmarks: A Cross-Benchmark Empirical Study.'' Research Square preprint, posted 18 August 2026. DOI: 10.21203/rs.3.rs-10665114/v1. CC BY 4.0; replication artifact \texttt{SQJReplicationArtifactv1.zip}. \url{https://doi.org/10.21203/rs.3.rs-10665114/v1}
\item Zhang, Xing, et al. ``Who Grades the Grader? Co-Evolving Evaluation Metrics and Skills for Self-Improving LLM Agents.'' arXiv preprint arXiv:2607.12790 (2026). \url{https://arxiv.org/abs/2607.12790}
\item Liu, Sibo. ``Audience-Bound Persistent Memory: Authorization Across the Memory Lifecycle.'' arXiv preprint arXiv:2609.36373 (2026). \url{https://arxiv.org/abs/2609.36373}
\item Cheng, Yezhou, Runjia Du, Zeming Liu, Qibai Chen, Hang Lyu, Yankai Zeng, Yilan Wei, and Bojun Lin. ``Scope Before You Persist: Preventing Cross-Family Interference in Agent Memory.'' arXiv preprint arXiv:2609.29144 (2026). \url{https://arxiv.org/abs/2609.29144}
\item Singh, Vivek Kumar, Preeti Priyam, and Gautam Bhowmick. ``ChurnBench: A Drift-Aware Benchmark Demonstrating That Refresh Scheduling, Not Cache Age, Governs Staleness in Agentic AI.'' arXiv preprint arXiv:2609.11515 (2026). \url{https://arxiv.org/abs/2609.11515}
\item Mughal, Ali Hassaan, and Muhammad Bilal. ``LLM-Based Test Oracles: Source-of-Authority Taxonomy --- A Systematic Literature Review.'' arXiv preprint arXiv:2607.05031 (2026). The arXiv \texttt{journal\_ref} names IEEE Access, vol. 14, 2026; the publisher DOI did not resolve during this audit, so the preprint identifier is the citation of record. \url{https://arxiv.org/abs/2607.05031}
\item Chitan, Florin Adrian. ``Session Risk Memory (SRM): Temporal Authorization for Deterministic Pre-Execution Safety Gates.'' arXiv preprint arXiv:2603.22350 (2026). \url{https://arxiv.org/abs/2603.22350}
\item Ahad, Tanzim, Ismail Hossain, Md Jahangir Alam, Sai Puppala, Syed Bahauddin Alam, and Sajedul Talukder. ``The Misattribution Gap: When Memory Poisoning Looks Like Model Failure in Agentic AI Systems.'' arXiv preprint arXiv:2605.22842 (2026). \url{https://arxiv.org/abs/2605.22842}
\item Hamidi, Asma, Michael Konstantinou, Renzo Degiovanni, and Mike Papadakis. ``How effective are traditional test criteria at detecting bugs in large language models generated code?'' arXiv preprint arXiv:2609.09315 (2026). \url{https://arxiv.org/abs/2609.09315}
\item Mishra, Shweta. ``Boundary-Mutation Testing for Pattern-Based Secret Detection: A Rule-Level Method and Cross-Scanner Evaluation.'' arXiv preprint arXiv:2609.02983 (2026). \url{https://arxiv.org/abs/2609.02983}
\item Girda, Nora, and Adrian Groza. ``Review Before Trust: Source-Grounded Integrity Gates for AI-Assisted Personal Health Records.'' arXiv preprint arXiv:2608.29965 (2026). \url{https://arxiv.org/abs/2608.29965}
\item de Oliveira, Rodrigo, Federico Pittino, James Gwinnutt, and Jay Nanavati. ``Beyond Correctness: Validity-Oriented Evaluation of Biomedical LLM Judges.'' arXiv preprint arXiv:2608.29127 (2026). \url{https://arxiv.org/abs/2608.29127}
\item Cho, Steven, Stefano Ruberto, and Valerio Terragni. ``Metamorphic Testing of Large Language Models for Natural Language Processing.'' Proceedings of the 2025 IEEE International Conference on Software Maintenance and Evolution (ICSME), pp. 174--186 (2025). DOI: 10.1109/ICSME64153.2025.00025. \url{https://arxiv.org/abs/2511.02108}
\end{enumerate}
\appendix
\section{History of the two control re-issues}
\label{app:reissue-history}
Section 4 reports what the target-alignment control now measures. This appendix reports how it got there, because the sequence is the clearest instance in this study of the register's recursive entry: an instrument that reports the absence of vacuity is subject to the same failure modes it reports on. An independently published verifier audit reports the same shape---a tool built to find checks that pass without checking itself produced such a check, and the failure surfaced as the best score the instrument could report (Hill, 2026, \S7.4). The wording here paraphrases that finding rather than quoting it, because the quotation could not be re-verified against the source by this audit; the section heading and the paper's definition of vacuous pass were verified, and the definition is quoted in Section 2.

The first re-issue found that its predecessor's headline conjunct compared a context digest against itself. Both cases in the comparison loaded the same prediction, so the equality held by construction and no input to the audited system could have broken it. The conjunct was demoted from a check to a disclosure and replaced by a paired form.

The second re-issue found a further conjunct half of the same shape: the stale branch is an unmutated deep copy, so its field difference against the frozen record is empty by construction. By then this manuscript had begun to claim that no conjunct survived as a theorem, which made that claim itself an assertion no perturbation reached. Finding it required perturbing the completeness claim rather than only the checks beneath it, and the manuscript had to be re-issued to carry the correction. Two re-issues exist and no more, so the ordinals are checkable against the two script versions retained on disk.

An earlier revision of this history said ``third'' in one paragraph and ``second'' in another, and the compiled PDF carried both, so a reader could not tell which the authors believed. That is recorded here rather than removed: a document that silently corrects its own ordinals leaves no way to check that the correction was the right one.

The general lesson is the recursive entry itself. A completeness claim needs its own perturbation battery, and the only defence found was to treat the claim that no vacuous conjunct remained as itself a conjunct to be attacked.

\section{Artifact field index}
\label{app:availability-index}
Every count quoted in the body resolves to a field named below. File names and field names are reproduced verbatim and may carry legacy internal identifiers such as \texttt{slots}, \texttt{pilot} or \texttt{twin}; the prose consistently says nested arms and development batch. Percentages and reductions computed from these fields are marked as derived where they appear in the body.
\begin{sloppypar}
The draft is backed by the local experiment bundle \texttt{e0\_observation}. The artifact counts quoted in the text resolve to the named fields below; percentages and reductions computed from them are marked as derived in the prose: \path{authoring_pilot_acoustic_v1/validation_report.json} (\texttt{projection\_conditions}, 144); \path{twin_v4/twin_gate_report_v4.json} (\texttt{status} \texttt{PASS\_REFERENCE\_ONLY}, \texttt{perturbation\_sensitive\_arms} 216, \texttt{families\_with\_detection\_power} 9, \texttt{visible\_positive\_controls} 9), with \path{twin_v4/twin_gate_detail_v4.json} supplying the per-family \texttt{hidden\_state\_windows} that reduce 216 to 9 distinct perturbations, and every \texttt{detection\_power} entry computed at run time rather than asserted; it supersedes \path{twin_v3/twin_gate_report_v3.json}, whose \texttt{status} was the unqualified \texttt{PASS} that Table~\ref{tab:gate-status} does not carry; \path{validity_v1/validity_gate_report_v1.json} (\texttt{arms\_tested} 72, \texttt{attested\_sensitive\_arms} 62, \texttt{attested\_reaction\_required\_arms} 50, \texttt{hidden\_twin\_comparisons} 24), with \path{validity_v1/validity_gate_detail_v1.json} supplying \texttt{hidden\_power\_arms} 10/0/0; \path{validity_v3/validity_gate_gap_analysis_v2.json} (\texttt{status} REPRODUCED, \texttt{self\_test\_status} \texttt{SELF\_TEST\_PASS}, \texttt{arm\_attribution} for the 12, \texttt{hidden\_comparisons} and \texttt{hidden\_comparison\_denominator} for the 10 informative of 24, \texttt{power\_coverage}, \texttt{per\_family\_hidden\_power}, \texttt{source\_predicates}, and \texttt{self\_validation} against every frozen aggregate), which supersedes \path{validity_v2/validity_gate_gap_analysis_v1.json}; the collinearity statement in Section 4 comes from \texttt{hidden\_comparison\_denominator} in this artifact and from nowhere else, and the artifact records that the resulting figure is not quotable as a detection rate; \path{results/gatemem_rag_policy_stub_household/summary.json} (\texttt{n\_checkpoints} 552, \texttt{n\_privacy}, \texttt{n\_safety} and \texttt{n\_utility} 184 each, \texttt{privacy\_context\_leakage\_rate}, \texttt{deletion\_context\_leakage\_rate}, and both answer-leakage-rate fields at 0.0), which supplies every replay count and rate in Section 4 and which an earlier revision of this draft quoted without naming; \path{sealed_timeline_v1/recipient_gate_report_v3.json} (\texttt{isolation\_checks} 27); \path{e2_sealed_full_v1_response_usability.json} (\texttt{usable} 217, \texttt{empty} 68, \texttt{nonempty\_truncated} 3, \texttt{truncated} 71, \texttt{slots} 288, \texttt{model}); \path{v1_5_admitted_validation_report.json} (\texttt{projection\_conditions} 216 for the coverage layer); \path{peer_review_revision_v3/gatemem_target_alignment_report_v3.json} (\texttt{control\_checks} 13, \texttt{mutation\_count} 15, \texttt{vacuous\_mutations} empty, \texttt{equivalence\_operator\_count} 6, \texttt{broken\_equivalence\_count} 3, \texttt{boundary\_guard\_class\_extracted\_from\_source}, \texttt{retracted\_assertion}, \texttt{manifest\_drift}), which supersedes \path{peer_review_revision_v2/gatemem_target_alignment_report_v2.json} and \path{peer_review_revision_v1/gatemem_target_alignment_report.json}; \path{peer_review_revision_v3/confirmatory_query_readiness_report_v3.json} (\texttt{status} FAIL, \texttt{query\_count} 84, \texttt{placeholder\_query\_count} 84, \texttt{family\_count} 24, \texttt{family\_floor}, \texttt{min\_queries\_per\_family}, \texttt{leakage\_token\_count} 0, \texttt{detection\_power.verdict} PARTIAL, \texttt{inconclusive\_families}, \texttt{status\_checks}, \texttt{self\_test\_status} \texttt{SELF\_TEST\_PASS}), which supersedes \path{confirmatory_query_readiness_report_v2.json} and is the readiness state the next phase would inherit: it forces a non-PASS verdict while any confirmatory query remains a placeholder, so the FAIL above is a property of the gate rather than of this draft's optimism. Every count in the corpus paragraph of Section 4 resolves to \path{corpus_matcher_v1/gatemem_corpus_matcher_report_v3.json} (\texttt{denominators}, \texttt{suppression}, \texttt{separator\_variants}, \texttt{not\_adjudicated}, \texttt{boundary\_guard\_class\_extracted\_from\_source}, \texttt{upstream\_intent\_evidence}, \texttt{corpus\_sha256}, \texttt{metrics\_sha256}, \texttt{gatemem\_commit}); its per-instance file \path{corpus_matcher_v1/gatemem_corpus_matcher_instances_v3.jsonl} carries upstream target values and checkpoint identifiers and is therefore \emph{not} redistributed, being reproducible from the pinned upstream commit by the named script. The zero verdict-flip figure resolves to \path{verdict_impact_v1/gatemem_verdict_impact_report_v1.json} (\texttt{status} \texttt{MEASURED\_ATTRIBUTION\_COMPLETE}, \texttt{pattern\_sets}, \texttt{answer\_side\_excluded}, \texttt{source\_grounding}, and per-arm \texttt{attribution} and \texttt{verdicts\_that\_would\_flip}), which derives its pattern sets from the upstream scorer's own source and refuses to run when that derivation fails; a sidecar \path{corpus_matcher_v1/gatemem_corpus_matcher_output_binding_v1.json} records that the corpus matcher's report does not carry its output file's digest, so the binding is stated externally rather than left implicit. \path{GATE_STATUS_AND_KNOWN_BAD_REGISTER_2026-10-05.md} records the full digest of every gate script and the status of every verdict cited in this draft. Artifact file names and field names above are reproduced verbatim and may carry legacy internal identifiers such as \texttt{slots}, \texttt{pilot}, or \texttt{twin}; the prose consistently says nested arms and development batch. One identifier class is worse than legacy: a character name from the host project appears four times in each language version of this draft---three times as family identifiers in Table~\ref{tab:family-usability} and once inside the artifact path \path{v1_5_admitted_sub-538bf0c712}---so this draft's own text carries identifiers that the de-identification condition on any release of it would have to strip. The denylist also matches one real place name per language version, which the transformation reports but does not alter, so the current derived copies remain 0 ADMIT and 2 EXCLUDE before a release decision. Publication therefore requires an opaque-ID transformation on derived copies, a human release decision for that retained class, and the pre-decision scanner result; the divergence from the frozen originals must be recorded, and the originals must not be rewritten. The scope of that transformation should not be overstated: it strips canon character identifiers from the research record, and it does not make the host project unidentifiable, because the accompanying public repository's own history names that project and because the real place names in this draft are retained by design. \emph{De-identified} here is therefore a statement about the character identifiers, not about the host project. That a document can state a release condition and violate it in the same paragraph is the failure mode this paper audits, so it is disclosed here rather than fixed silently. The frozen allocation and artifact manifest are the authoritative records for this stage. The accompanying literature recheck \path{novelty_recheck_2026-10-04.md} records the searches behind Section 2, including the channels that were rate-limited. That search is non-exhaustive: the protocol record \path{novelty_query_protocol_v1_2026-10-05.md} requires atomic queries, one known-occupier calibration query per round, recorded result totals, and scope/claims reading for every surviving neighbour. The external GateMem code snapshot is MIT-licensed and its dataset is marked CC-BY-4.0 by the upstream repository and dataset card; upstream benchmark files, complete answers and hidden checkpoints are not redistributed in this artifact. Production secrets and runtime world ledgers are not inputs.
\end{sloppypar}

\section{Family-level response usability}
\label{app:family-usability}
Descriptive only. Each row contains 24 nested arms while the intended independent unit remains the family. No confidence interval or significance test is reported at the arm level, because arms are nested within families and arm-level inference would be pseudoreplicated (Hurlbert, 1984). Nothing in this table is interpreted as causal, and none of it may be before semantic rating exists.

\begin{table}[h]
\centering
\caption{Family-level response usability (descriptive).}
\label{tab:family-usability}
\scriptsize
\begin{tabular}{lrrrrr}
\toprule
Family & Nested arms & Usable & Empty & Length finish & Usable rate \\
\midrule
\texttt{access-f01} & 24 & 18 & 5 & 6 & 75.00\% \\
\texttt{access-f02} & 24 & 13 & 11 & 11 & 54.17\% \\
\texttt{access-f03} & 24 & 15 & 9 & 9 & 62.50\% \\
\texttt{communication-synthetic-01} & 24 & 24 & 0 & 0 & 100.00\% \\
\texttt{communication-synthetic-02} & 24 & 23 & 1 & 1 & 95.83\% \\
\texttt{communication-synthetic-03} & 24 & 23 & 1 & 1 & 95.83\% \\
\texttt{coverage-f01} & 24 & 15 & 9 & 9 & 62.50\% \\
\texttt{coverage-f02} & 24 & 15 & 9 & 9 & 62.50\% \\
\texttt{coverage-f03} & 24 & 17 & 7 & 7 & 70.83\% \\
\texttt{crossday-f01} & 24 & 18 & 4 & 6 & 75.00\% \\
\texttt{crossday-f02} & 24 & 17 & 7 & 7 & 70.83\% \\
\texttt{crossday-f03} & 24 & 19 & 5 & 5 & 79.17\% \\
\bottomrule
\end{tabular}
\end{table}

\end{document}
